\section{Introduction}
\subsection{Two matrices attached to a hyperplane arrangement}
Let $\C=\{H_i\}_{i\in E_0}$ be an essential arrangement of distinct
affine hyperplanes in $\R^r$. We denote its chambers by $\cF$ and its
bounded chambers by $\cB\subseteq\cF$.

For a finite polytopal complex $F$, the $h$-polynomial of $F$ is defined as
\[
 h(F,t)=\sum_{\varnothing\ne G\leq F}(t-1)^{\dim G},
 \qquad h(\varnothing,t)=0,
\]
where $G$ runs over its nonempty faces. For a closed face $F$, let
$F_\bd$ be its complex of bounded faces. For chambers
$P,Q\in\cF$, let $\Sep(P,Q)$ be the set of hyperplanes separating
them and put $d(P,Q)=|\Sep(P,Q)|$. 

Consider the polynomial rings $\ZM[q]$ and $\ZM[\mathbf b] = \ZM[b_i~|~i\in E_0]$ in variables $q$ and $b_i$ for $i \in E_0$ respectively.   As shorthand, for a subset $J \subseteq E_0$, we write $b_J=\prod_{i\in J}b_i$.

Our main results involve two $\cF\times \cF$-matrices. 

\begin{definition}\cite{Varchenko1993} The \textit{Varchenko matrix} $V_b$ is the $\cF\times \cF$-matrix with $(P,Q)$-entry
\[
 V_b(P,Q)=b_{\Sep(P,Q)} \in \ZM[\mathbf b],
\]
for any pair of chambers $P,Q\in\cF$.  Let $V_q=V_b|_{b_i=q}$ be the matrix with entries in $\ZM[q]$ obtained from $V_b$ by specializing all of the variables $b_i$ to $q$.
\end{definition}

\begin{definition} The \textit{$q$-intersection matrix} $W_q$ is the $\cF\times \cF$-matrix with $(P,Q)$-entry
\[
W_q(P,Q)=(-q)^{d(P,Q)}h((\bar P\cap\bar Q)_\bd,q^2) \in \ZM[q],
\]
for any pair of chambers $P,Q\in\cF$.  
\end{definition}

For an $\cF \times \cF$-matrix $X$, write $X^{\cB}$ for its
$\cB\times\cB$-submatrix. For hyperplane arrangements with normal crossings, the $\cB \times \cB$-matrix $S_q:=W_q^{\cB}$ is the $q$-intersection matrix introduced previously in \cite{EM}.

Our main theorems are a determinant formula for $(V_b^{-1})^{\cB}$, the $\cB \times \cB$-submatrix of the inverse of the Varchenko matrix, confirming a conjecture of Lam \cite[Conjecture 8.10]{Lam}, and a determinant formula for $S_q$ for normal crossing arrangements, confirming our conjecture~\cite[Conjecture 1.7]{EM}.

\subsection{Determinant Formulas}
Denote by $\widehat{\C}$ and $\C_0$ the homogenized and central
arrangements attached to $\C$, and by $M$ and $M_0$ the matroids of
their normal vectors. Their ground sets are $E=E_0\sqcup\{0\}$ and
$E_0$, respectively, where $0$ labels the hyperplane at infinity.
Write $\cL = \cL(M)$ for the lattice of flats of $M$, $\cL_0 \subseteq \cL$ for the set of flats containing $0$,
and $\Z(M_0)$ for the cyclic flats of $M_0$.

For a matroid $N$ we write $\beta(N)$ for Crapo's beta invariant and $\mup(N)$ for the unsigned Möbius number.  For any flat $F$ of $N$, we set \[w_F=\beta(N^F)\beta(N_F),\] where $N^F=N|F$ denotes the restriction to $F$ and $N_F=N/F$ denotes the contraction by $F$.  We recall additional matroid notation in Section~\ref{sec:background}.

The following theorem proves Lam's determinant conjecture~\cite[Conjecture 8.10]{Lam} in the case of real affine hyperplane
arrangements.
\begin{theorem}\label{thm:Lam-real}
For any essential real affine hyperplane arrangement,
\begin{equation}\label{eq:Lam-determinant}
 \det\bigl((V_b^{-1})^{\cB}\bigr)
 =\frac{\displaystyle\prod_{F\in\cL_0\setminus\{E\}}
          (1-b_{E\setminus F}^2)^{w_F}}
        {\displaystyle\prod_{F\in\cL\setminus(\cL_0\cup\{\varnothing\})}
          (1-b_F^2)^{w_F}}.
\end{equation}
\end{theorem}

Lam states his conjecture in terms of the bounded block of the Betti homology matrix, which he proves is equal to $(-1)^r(V_b^{-1})^{\cB}$, as recalled in
Section~\ref{subsec:betti-pairing}.

We also prove the determinant formula conjectured in
\cite[Conjecture 1.7]{EM}.  To state our result, we write
\[
 \qint{n}=1+q^2+\cdots+q^{2n-2}=\frac{1-q^{2n}}{1-q^2}.
\]


\begin{theorem}\label{thm:intersection-real}
If $\C$ is an essential real affine hyperplane arrangement with normal crossings, that is, if any non-empty intersection of $s$ hyperplanes in $\C$ has codimension $s$, then
\begin{equation}\label{eq:Sq-determinant}
 \det S_q
 =\prod_{\substack{K\in\Z(M_0)\\K\ne E_0}}
   \qint{|E_0\setminus K|}^{\beta(M_{0,K})\mup((M_0^K)^*)}.
\end{equation}
\end{theorem}

\begin{remark}
Specializing $q$ to 1, the entries of $S_1=S_q|_{q=1}$ are, up to sign, counts of vertices of common
faces. The $q=1$ determinant formula was proved in \cite[Theorem 1.4]{EM}
by identifying $S_1$ with the Brylawski--Schechtman--Varchenko form~\cite{SchechtmanVarchenko1991,BrylawskiVarchenko1997}
for the dual matroid $M_0^*$, with all weights equal to one.
For a unimodular arrangement, the bounded chambers index the toric
components of the compact core of a smooth hypertoric variety.
The matrix $S_1$ describes the intersection form on the top cohomology of this core
\cite{HauselSwartz2006}.

For realizable (and more generally orientable) matroids, the matrix $S_1$ can be used to study the representation theory of the matroidal Schur algebras of~\cite{Mautner,BradenMautner}.  We expect the existence of matroidal $q$-Schur algebras for which $S_q$ should play an analogous role.
\end{remark}


The proof of Theorem~\ref{thm:Lam-real} depends on a determinant formula for a de Rham pairing and many other results in the monumental work of Lam \cite{Lam} as well as the beautiful formulas for the determinants of period matrices due to Varchenko, Loeser--Sabbah and Douai--Terao \cite{Varchenko1990,LoeserSabbah1991,DouaiTerao1997}. We deduce Theorem~\ref{thm:intersection-real} by relating, for arrangements with normal crossings, the $q$-intersection matrix $W_q$ to the inverse of the Varchenko matrix, as we describe below.

\begin{remark}
Both determinant formulas are conjectured to hold in the more general setting of affine oriented
matroids, with an analogous normal crossing hypothesis for the second (see Definition~\ref{def:normal-crossing}). Our proofs
in the setting of hyperplane arrangements use realizability through the period determinant formula.
\end{remark}


\subsection{Koszul duality and the inverse Varchenko matrix}  To show that Theorem~\ref{thm:Lam-real} implies Theorem~\ref{thm:intersection-real}, we prove that when $\C$ has normal crossings, then the $q$-Varchenko matrix $V_q$ and the $q$-intersection matrix $W_q$
are inverse up to the scalar factor $(1-q^2)^r$.
In fact, we show the following stronger, categorical statement. Namely, we define the \emph{taut path algebra}
$\Taut$ and the \emph{intersection algebra} $\Inter$, and show that their Hilbert
matrices are $V_q$ and $(1-q^2)^{-r}W_{-q}$, respectively.
These algebras are Koszul dual by Theorem~\ref{thm:chamber-koszul}.
The Hilbert-series identity for Koszul dual algebras then gives the
following comparison.
\begin{theorem}\label{thm:intro-inverse}
For a normal crossing arrangement $\C$, we have
\[
 W_q=(1-q^2)^rV_q^{-1},\qquad
 S_q=(1-q^2)^r(V_q^{-1})^{\cB}.
\]
\end{theorem}
We prove this comparison more generally for normal crossing affine oriented matroids (see Definition~\ref{def:normal-crossing}).
In this setting as well, it reduces the conjectural determinant formula
\eqref{eq:Sq-determinant} to \eqref{eq:Lam-determinant} in the normal
crossing case.

\begin{remark}
The algebras $\Taut$ and $\Inter$ are closely related to algebras introduced by Braden--Licata--Proudfoot--Webster
\cite{BLPWgale,BradenLicataProudfootWebster2012}. In the realizable
normal crossing case, $\Inter$ appears in~\cite[Section 4.5]{BLPWgale} as the extended equivariant algebra
$\widetilde B_{\mathrm{ext}}$ defined on all feasible chambers.  The Koszul duality of Braden--Licata--Proudfoot--Webster concerns category $\mathcal O$ for Gale dual polarized arrangements.
Our theorem gives a duality on the full chamber set of a single affine
oriented matroid. For real arrangements, $S_q$ is also the graded
Euler form on the simple objects indexed by bounded chambers in
hypertoric category $\mathcal O$ \cite[Theorem 3.1]{EM}.
\end{remark}

\subsection{Outlook}\label{subsec:outlook}

\begin{enumerate*}[label=(\arabic*), mode=unboxed, itemjoin=\endgraf\noindent]
\item
It would be interesting to extend the Koszul duality to hyperplane arrangements (and affine oriented matroids) that do not have normal
crossings. The natural replacement for the intersection algebra is
\[
 \operatorname{Ext}_{\Taut}^{\bullet}(\Taut_0,\Taut_0),
\]
together with its internal multigrading and higher multiplications.
Taking alternating sums in cohomological degree recovers $V_b^{-1}$,
thereby categorifying Lam's Betti homology matrix at $b_0=0$, up to
the sign $(-1)^r$. For the arrangements describing $M_{0,n}$, the
bounded-chamber entries are string-corrected biadjoint scalar amplitudes
up to normalization \cite[Section 19.4]{Lam}, so this also gives a
categorification of these amplitudes. An explicit combinatorial
description and an extension to general $b_0$ would be desirable.

\item
A combinatorial proof avoiding the period determinant formula of
Loeser--Sabbah and Douai--Terao could establish both determinant formulas for
affine oriented matroids, with the normal crossing assumption for
the $q$-intersection formula. We already verified the latter formula
for the non-realizable pseudoline arrangement pictured on the cover
of \emph{Oriented Matroids}
\cite{OrientedMatroids1999}, and for a generic affine lifting of an
orientation of the V\'amos matroid. See \cite[Sections A.2--A.3]{EM}.

\item
It would also be interesting to apply our methods to the Koszulity
question raised by Kowalenko--Mautner for the algebras associated to
generic oriented matroid programs \cite[Section 1.3]{KowalenkoMautner}.
They establish quadratic duality in general and Koszulity for
Euclidean programs. An affirmative answer would extend the Koszul
duality of Braden--Licata--Proudfoot--Webster \cite{BLPWgale} to arbitrary
generic oriented matroid programs.
\end{enumerate*}

\subsection{Structure}
Section~\ref{sec:background} fixes notation for arrangements and affine
oriented matroids. Section~\ref{sec:periods} recalls the twisted
intersection forms and period determinant formula used in
Section~\ref{sec:determinants} to prove Theorem~\ref{thm:Lam-real}.
Section~\ref{sec:comb} collects the combinatorial lemmas for normal
crossing affine oriented matroids. Section~\ref{sec:koszul} proves
Koszul duality and the inverse-matrix identity.
Section~\ref{sec:qdet} deduces Theorem~\ref{thm:intersection-real}.
